% !PRINTED 78
\weberpage{78}
%<<CONT
considered as functions of $u$, related $T$-functions of the second order with the same periodicity factors, and hence any three of them are linearly dependent. We therefore have
\[
A\,\vartheta_{00}(u)\,\vartheta_{00}(u+v) + B\,\vartheta_{10}(u)\,\vartheta_{10}(u+v) + C\,\vartheta_{11}(u)\,\vartheta_{11}(u+v) = 0,
\]
from which, for $u=0$, $u=\dfrac{1}{2}$:
\[
\begin{aligned}
A\,\vartheta_{00}\vartheta_{00}(v) + B\,\vartheta_{10}\vartheta_{10}(v) &= 0,\\
A\,\vartheta_{01}\vartheta_{01}(v) + C\,\vartheta_{10}\vartheta_{10}(v) &= 0
\end{aligned}
\]
and from these:

\weq{11}{\begin{aligned}\vartheta_{10}\vartheta_{10}(v)\,\vartheta_{00}(u)\,\vartheta_{00}(u+v) - \vartheta_{00}\vartheta_{00}(v)\,\vartheta_{10}(u)\,\vartheta_{10}(u+v)\\ {}- \vartheta_{01}\vartheta_{01}(v)\,\vartheta_{11}(u)\,\vartheta_{11}(u+v) = 0.\end{aligned}}

In the same way, if $w$ is a third variable, the products contained in the form
\[
\vartheta_{g_1,g_2}(u+v)\,\vartheta_{g_1,g_2}(u+w), \quad \vartheta_{g_1,g_2}(u)\,\vartheta_{g_1,g_2}(u+v+w)
\]
are, considered as functions of $u$, related $T$-functions of the second order with the same periodicity factors, so that between any three of them there exists a linear dependence. The coefficients are determined, as above, by special values of $u$. Thus the system of formulae follows:

\weq{12}{\begin{aligned}\vartheta_{00}\vartheta_{00}(v+w)\,\vartheta_{00}(u+v)\,\vartheta_{00}(u+w)\\ {}= \vartheta_{00}(u)\vartheta_{00}(v)\vartheta_{00}(w)\vartheta_{00}(u+v+w)\\ {}- \vartheta_{11}(u)\vartheta_{11}(v)\vartheta_{11}(w)\vartheta_{11}(u+v+w).\end{aligned}}
\weq{13}{\begin{aligned}\vartheta_{01}\vartheta_{01}(v+w)\,\vartheta_{01}(u+v)\,\vartheta_{01}(u+w)\\ {}= \vartheta_{01}(u)\vartheta_{01}(v)\vartheta_{01}(w)\vartheta_{01}(u+v+w)\\ {}+ \vartheta_{11}(u)\vartheta_{11}(v)\vartheta_{11}(w)\vartheta_{11}(u+v+w).\end{aligned}}
\weq{14}{\begin{aligned}\vartheta_{10}\vartheta_{10}(v+w)\,\vartheta_{10}(u+v)\,\vartheta_{10}(u+w)\\ {}= \vartheta_{10}(u)\vartheta_{10}(v)\vartheta_{10}(w)\vartheta_{10}(u+v+w)\\ {}+ \vartheta_{11}(u)\vartheta_{11}(v)\vartheta_{11}(w)\vartheta_{11}(u+v+w).\end{aligned}}

All these formulae can be regarded as special cases of a general formula which \gesperrt{Jacobi} derived from the series developments and employed as the foundation of the theory of elliptic functions\footnote{Theorie der elliptischen Funktionen, aus den Eigenschaften der Thetareihen abgeleitet. \gesperrt{Jacobis} gesammelte Werke, Bd.\,I, S.\,497.}. This formula can also be obtained in the following way from the concept of the theta functions.
